% !TeX program = xelatex
% 第一二问依据冻结源码及结果核对；附录源程序段保持原文。
\documentclass[a4paper]{article}
\usepackage{geometry}
\geometry{a4paper,left=2.7cm,right=2.7cm,top=2.54cm,bottom=2.54cm}
 

\usepackage[UTF8,heading=true]{ctex}
 

\usepackage{abstract}
\setlength{\abstitleskip}{0em}
\setlength{\absleftindent}{0pt}
\setlength{\absrightindent}{0pt}
\setlength{\absparsep}{0em}
\renewcommand{\abstractname}{\textbf{\zihao{4}{摘要}}}
\renewcommand{\abstracttextfont}{\zihao{-4}} %设置摘要正文字号
 
\usepackage{fancyhdr}
\pagestyle{plain}
 

\usepackage{amssymb}
\usepackage{amsmath}
 

\usepackage{float}
\usepackage{subfig}
\captionsetup[figure]{labelsep=space} %去除图标题的冒号
\captionsetup[table]{labelsep=space} %去除表格标题的冒号
%设置标题格式
\ctexset {
	%设置一级标题的格式
	section = {
		name={,、},
		number=\chinese{section}, %设置中文版的标题
		aftername=,
	},
	%设置三级标题的格式
	subsubsection = {
		format += \zihao{-4} % 设置三级标题的字号
	}
}
 
 
%使得英文字体都为Time NewTown
%\usepackage{times}
 
%图片包导入
\usepackage{graphicx}
\graphicspath{{figures/}}
% 有原图目录时优先使用原文件；否则使用同目录 paper_figures.pdf 的对应页。
\newcommand{\paperfigure}[3]{%
\IfFileExists{figures/#2}{\includegraphics[width=#1\textwidth]{figures/#2}}%
{\includegraphics[width=#1\textwidth,page=#3]{paper_figures.pdf}}}
 %图片在当前目录下的figures目录
 
%参考文献包引入
\usepackage[numbers,sort&compress]{natbib}
 
%代码格式
\usepackage{listings}
\usepackage{xurl}
\usepackage{xcolor}%代码配色，无额外高亮依赖
\lstset{
	%numbers=left, %设置行号位置
	%	numberstyle=\tiny, %设置行号大小
	keywordstyle=\color{blue}, %设置关键字颜色
	commentstyle=\color[cmyk]{1,0,1,0}, %设置注释颜色
	escapeinside=``, %逃逸字符(1左面的键)，用于显示中文
	breaklines, %自动折行
	extendedchars=false, %解决代码跨页时，章节标题，页眉等汉字不显示的问题
	xleftmargin=1em,xrightmargin=1em, aboveskip=1em, %设置边距
	tabsize=4, %设置tab空格数
	showspaces=false %不显示空格
}
 
 

\lstdefinestyle{appendixcode}{
    language=Python,
    basicstyle=\ttfamily\fontsize{5.2}{5.8}\selectfont,
    numbers=left,
    numberstyle=\fontsize{4.5}{5}\selectfont,
    numbersep=5pt,
    stepnumber=1,
    firstnumber=1,
    breaklines=true,
    breakatwhitespace=false,
    columns=fullflexible,
    keepspaces=true,
    showstringspaces=false,
    tabsize=4,
    frame=none,
    xleftmargin=1.4em,
    xrightmargin=0.1em,
    aboveskip=0.35em,
    belowskip=0.55em,
    escapeinside={(*@}{@*)}
}

\renewcommand{\refname}{}
 
%item包
\usepackage{enumitem}
 
%表格加粗
\usepackage{booktabs}
 
%设置表格间距
\usepackage{caption}
 
%允许表格长跨页
\usepackage{longtable}
\usepackage{array}
\usepackage{needspace}
\raggedbottom
\setlength{\emergencystretch}{2em}
\IfFontExistsTF{DejaVu Sans Mono}
  {\setmonofont{DejaVu Sans Mono}[Scale=MatchLowercase]}
  {\setmonofont{Latin Modern Mono}}
% 第三问原正文使用的数学宏；只补定义，不修改公式含义。
\providecommand{\measHist}{\mathcal D}
\providecommand{\meas}{\widehat{\theta}}
\providecommand{\effErr}{\varepsilon_{\mathrm{eff}}}
\providecommand{\cross}{\operatorname{cross}}
\providecommand{\locRegion}{P}
\providecommand{\emptyHist}{\mathcal Z}
\providecommand{\proj}{\widetilde P}
\providecommand{\knownSet}{\mathcal K}
\providecommand{\clearedSet}{\mathcal L}
\providecommand{\emptySet}{\mathcal A}
\providecommand{\chan}{\mathcal H}

 
%伪代码用到的宏包
\usepackage{algorithmic}
\usepackage{algorithm}
 

% 图文整合版：正文仍为小四，调整空白，不缩小源码正文或更改代码。
\linespread{1.12}
\ctexset{
  section={beforeskip=1.5ex plus .3ex minus .2ex,afterskip=.8ex plus .2ex},
  subsection={beforeskip=1.2ex plus .2ex,afterskip=.5ex plus .1ex},
  subsubsection={beforeskip=.9ex plus .2ex,afterskip=.4ex plus .1ex}}
\setlist{itemsep=1pt,topsep=3pt,parsep=0pt,partopsep=0pt}
\setlength{\textfloatsep}{9pt plus 2pt minus 2pt}
\setlength{\floatsep}{8pt plus 2pt minus 2pt}
\setlength{\intextsep}{8pt plus 2pt minus 2pt}
\renewcommand{\topfraction}{.90}
\renewcommand{\bottomfraction}{.80}
\renewcommand{\textfraction}{.08}
\renewcommand{\floatpagefraction}{.70}
\captionsetup{font=small,skip=4pt}
\AtBeginDocument{\setlength{\abovedisplayskip}{6pt plus 2pt minus 2pt}
\setlength{\belowdisplayskip}{6pt plus 2pt minus 2pt}
\setlength{\abovedisplayshortskip}{3pt plus 1pt}
\setlength{\belowdisplayshortskip}{5pt plus 1pt}}

%正文区
%<paper-block id="title.h001" type="heading">
\title{有界测向误差下干扰源的交会定位与搜索清除策略} 
%</paper-block>
\author{}
\date{} %不显示日期
 
%文档
